\documentclass[../../main.tex]{subfiles}
\graphicspath{{../../image/}{../../figure/}}
\begin{document}

\section{等熵弹性方程组}\label{section:ch17-等熵弹性}

当方程组\eqref{align:ch17-1731}为线性时，\cref{theorem:ch17-两律真非线性补偿紧致}的断言显然失效。另一方面，真非线性是过强的限制：它可允许在状态空间中沿有限条曲线失效，只要这些曲线横截相交于 Riemann 不变量的等值线。这一点将在一维等熵热弹性守恒律方程组\eqref{align:ch07-p系统}的框架下展示，
即
\begin{align}
\begin{cases}
\partial_t u-\partial_x v=0\\
\partial_t v-\partial_x\sigma(u)=0,
\end{cases}
\label{align:ch17-1771}
\end{align}
其假设为 $\sigma''(u)\neq0$（$u\neq0$），但 $\sigma''(0)=0$，以致真非线性在状态空间中沿直线 $u=0$ 失效。尽管如此，\cref{theorem:ch17-两律真非线性补偿紧致}的类比结果仍成立：

\begin{theorem}[等熵弹性的补偿紧致性]\label{theorem:ch17-弹性补偿紧致}
设 $\Omega$ 是 $\mathbb{R}^2$ 的开子集，$\{(u_k,v_k)\}$ 是 $L^\infty(\Omega;\mathbb{R}^2)$ 中的有界序列，满足
\begin{align}
\partial_t\eta(u_k,v_k)+\partial_x q(u_k,v_k)\subset\text{紧集于 }W^{-1,2}_{\mathrm{loc}}(\Omega),
\label{align:ch17-1772}
\end{align}
其中 $(\eta,q)$ 是\eqref{align:ch17-1771}的任意熵--熵流对。则存在子列 $\{(u_j,v_j)\}$，它在 $\Omega$ 上几乎处处收敛。
\end{theorem}

\begin{proof}
如\cref{theorem:ch17-两律真非线性补偿紧致}的证明中那样，提取子列 $\{(u_j,v_j)\}$ 并识别相联系的 Young 测度族 $\nu_{x,t}$。固定 $(x,t)\in\Omega$，并相对于 Riemann 不变量监测 $(x,t)$ 处的 Young 测度 $\nu$：
\begin{align}
z=\int_0^u[\sigma'(\omega)]^{\frac{1}{2}}\,d\omega+v,\qquad w=-\int_0^u[\sigma'(\omega)]^{\frac{1}{2}}\,d\omega+v.
\label{align:ch17-1773}
\end{align}
我们需证明包含 $\nu$ 支撑的最小矩形 $\mathscr{R}=[z^-,z^+]\times[w^-,w^+]$ 坍缩为单点。

沿\cref{theorem:ch17-两律真非线性补偿紧致}证明中那些不依赖于方程组真非线性的步骤回溯，我们重新导出\eqref{align:ch17-1769}。论证的余下部分将依赖于 $\mathscr{R}$ 与真非线性失效的直线 $z=w$ 的相对位置。

先设直线 $z=w$ 不与 $\mathscr{R}$ 的右边相交，即 $z^+\notin[w^-,w^+]$。此时，\eqref{align:ch17-17610}仍有效，给出 $z^-=z^+$。于是 $\mathscr{R}$ 坍缩为 $[z^+]\times[w^-,w^+]$，由我们的假设它完全位于真非线性区域内，从而由熟知的论证得 $w^-=w^+$，验证了定理的断言。当直线 $z=w$ 不碰到 $\mathscr{R}$ 其他三边中任一边时，类似论证适用。

余下需考察直线 $z=w$ 与 $\mathscr{R}$ 四边都相交的情形，即 $z^-=w^-$ 且 $z^+=w^+$。即使在这一情形下，由\eqref{align:ch12-估计1}，
$q_k-\lambda\eta_k$ 沿 $[z^-]\times[w^-+\varepsilon,w^+]$ 与 $[z^+]\times[w^-,w^+-\varepsilon]$ 不变号，因此熟知的论证仍成立，给出 $z^-=z^+$，除非测度 $\nu^-$ 与 $\nu^+$ 分别集中于顶点 $(z^-,w^-)$ 与 $(z^+,w^+)$。当这种情况发生时，\eqref{align:ch17-1769}退化为
\begin{align}
q(z^-,w^-)-\lambda(z^-,w^-)\eta(z^-,w^-)=q(z^+,w^+)-\lambda(z^+,w^+)\eta(z^+,w^+).
\label{align:ch17-1774}
\end{align}
特别地，对平凡熵--熵流对 $(u,-v)$ 应用\eqref{align:ch17-1774}。在“西南”顶点，$u^-=0$、$v^-=z^-=w^-$；而在“东北”顶点，$u^+=0$、$v^+=z^+=w^+$。于是\eqref{align:ch17-1774}给出 $z^-=z^+=w^-=w^+$。证明完成。
\end{proof}

$\sigma(u)$ 的光滑性一般不能放松，因为例子表明：当 $\sigma''(u)$ 在 $u=0$ 处不连续时，上述命题的断言可能失效。

特别地，当弹性介质的响应如同一个“硬弹簧”，即 $\sigma$ 在 $u<0$ 时凹、在 $u>0$ 时凸时，\cref{theorem:ch17-弹性补偿紧致}适用：
\begin{align}
u\sigma''(u)>0,\qquad u\neq0.
\label{align:ch17-1775}
\end{align}
对该情形，可对由消失黏性方法构造的近似解确立 $L^\infty$ 界，即它们作为 Cauchy 问题
\begin{align}
\begin{cases}
\partial_t u-\partial_x v=\mu\,\partial_x^2u\\
\partial_t v-\partial_x\sigma(u)=\mu\,\partial_x^2v,
\end{cases}
\label{align:ch17-1776}
\end{align}
\begin{align}
(u(x,0),v(x,0))=(u_{0\mu}(x),v_{0\mu}(x)),\qquad-\infty<x<\infty
\label{align:ch17-1777}
\end{align}
的解。

\begin{theorem}[等熵弹性的不变区域]\label{theorem:ch17-弹性不变区域}
在假设\eqref{align:ch17-1775}下，对任意 $M>0$，由
\begin{align}
\mathscr{U}_M=\{(u,v):-M\leqslant z(u,v)\leqslant M,\ -M\leqslant w(u,v)\leqslant M\}
\label{align:ch17-1778}
\end{align}
定义的集合 $\mathscr{U}_M$（其中 $z$ 与 $w$ 是\eqref{align:ch17-1771}的 Riemann 不变量\eqref{align:ch17-1773}）是\eqref{align:ch17-1776}、\eqref{align:ch17-1777}解的一个（正）不变区域。
\end{theorem}

\begin{proof}
标准证明基于最大值原理。这里给出另一个证明，它依赖于熵，因而更贴近双曲理论的精神。它的优点是对\eqref{align:ch17-1776}解的光滑性要求较低；此外，它可容易地推广到任何其他在\eqref{align:ch17-1771}凸熵下耗散的近似格式。

对系统\eqref{align:ch17-1771}，确定熵--熵流对 $(\eta,q)$ 的方程\eqref{align:ch07-熵流对定义}化为
\begin{align}
\begin{cases}
q_u(u,v)=-\sigma'(u)\eta_v(u,v)\\
q_v(u,v)=-\eta_u(u,v).
\end{cases}
\label{align:ch17-1779}
\end{align}
注意\eqref{align:ch17-1779}允许解族
\begin{align}
\eta_m(u,v)=Y_m(u)\cosh(mv)-1,
\label{align:ch17-17710}
\end{align}
\begin{align}
q_m(u,v)=-\frac{1}{m}Y_m'(u)\sinh(mv),
\label{align:ch17-17711}
\end{align}
其中 $m=1,2,\cdots$，而 $Y_m$ 是常微分方程
\begin{align}
Y_m''(u)=m^2\sigma'(u)Y_m(u),\qquad-\infty<u<\infty,
\label{align:ch17-17712}
\end{align}
在初值
\begin{align}
Y_m(0)=1,\qquad Y_m'(0)=0
\label{align:ch17-17713}
\end{align}
下的解。

简单计算给出
\begin{align}
\eta_{muu}\eta_{mvv}-\eta_{muv}^2\geqslant m^2\bigl[m^2\sigma'Y_m^2-Y_m'^2\bigr].
\label{align:ch17-17714}
\end{align}
此外，由\eqref{align:ch17-17712}，
\begin{align}
\bigl[m^2\sigma'Y_m^2-Y_m'^2\bigr]'=m^2\sigma''Y_m^2.
\label{align:ch17-17715}
\end{align}

因此，\eqref{align:ch17-1775}蕴含\eqref{align:ch17-17714}右端为正，从而 $\eta_m(u,v)$ 是 $\mathbb{R}^2$ 上的凸函数。此外，$\eta_m(0,0)=0$ 且 $\eta_{mu}(0,0)=\eta_{mv}(0,0)=0$，以致 $\eta_m(u,v)$ 正定。

接着考察 $\eta_m(u,v)$ 当 $m\to\infty$ 时的渐近性。变量替换 $(u,Y_m)\mapsto(\xi,X_m)$：
\begin{align}
\xi=\int_0^u[\sigma'(\omega)]^{\frac{1}{2}}\,d\omega,
\label{align:ch17-17716}
\end{align}
\begin{align}
X_m=(\sigma')^{\frac{1}{4}}Y_m,
\label{align:ch17-17717}
\end{align}
将\eqref{align:ch17-17712}化为
\begin{align}
\ddot X_m=m^2X_m+\Bigl[\tfrac{1}{4}(\sigma')^{-2}\sigma'''-\tfrac{5}{16}(\sigma')^{-3}(\sigma'')^2\Bigr]X_m,
\label{align:ch17-17718}
\end{align}
其渐近性由常数变易公式导出：
\begin{align}
X_m(\xi)=\Bigl[\sigma'(0)^{\frac{1}{4}}+O\bigl(\tfrac{1}{m}\bigr)\Bigr]\cosh(m\xi),
\label{align:ch17-17719}
\end{align}
当 $m\to\infty$、且 $\xi$ 限制在任意固定有界区间内时成立。

结合\eqref{align:ch17-17710}、\eqref{align:ch17-17717}、\eqref{align:ch17-17719}、\eqref{align:ch17-17716}与\eqref{align:ch17-1773}，我们推出
\begin{align}
\lim_{m\to\infty}\eta_m(u,v)^{\frac{1}{m}}=
\begin{cases}
\exp[z(u,v)],&\text{若 }u>0,\ v>0,\\
\exp[w(u,v)],&\text{若 }u<0,\ v>0,\\
\exp[-w(u,v)],&\text{若 }u>0,\ v<0,\\
\exp[-z(u,v)],&\text{若 }u<0,\ v<0.
\end{cases}
\label{align:ch17-17720}
\end{align}

现考虑\eqref{align:ch17-1776}、\eqref{align:ch17-1777}的解 $(u_\mu,v_\mu)$，其中 $(u_{0\mu},v_{0\mu})\in L^2(-\infty,\infty)$ 且取值于\eqref{align:ch17-1778}定义的区域 $\mathscr{U}_M$。我们写出\eqref{align:ch17-1737}，取 $U_\mu=(u_\mu,v_\mu)$、$\eta=\eta_m$、$q=q_m$，并在条带 $(-\infty,\infty)\times[0,t]$ 上积分，得
\begin{align}
\int_{-\infty}^\infty\eta_m(u_\mu(x,t),v_\mu(x,t))\,dx\leqslant\int_{-\infty}^\infty\eta_m(u_{0\mu}(x),v_{0\mu}(x))\,dx.
\label{align:ch17-17721}
\end{align}
将\eqref{align:ch17-17721}取 $1/m$ 次幂，令 $m\to\infty$ 并使用\eqref{align:ch17-17720}，我们断言 $(u_\mu(\cdot,t),v_\mu(\cdot,t))$ 取值于区域 $\mathscr{U}_M$。证明完成。
\end{proof}

上述命题连同\cref{theorem:ch17-弹性补偿紧致}给出方程组\eqref{align:ch17-1771}的一个存在性定理，它摆脱了外加假设：

\begin{theorem}[等熵弹性的存在性]\label{theorem:ch17-弹性存在性}
设 $(u_\mu,v_\mu)$ 是上半平面上初值问题\eqref{align:ch17-1776}、\eqref{align:ch17-1777}的解，其中 $(u_{0\mu},v_{0\mu})\rightharpoonup(u_0,v_0)$ 在 $L^\infty(-\infty,\infty)$ 弱$^*$ 意义下成立。在条件\eqref{align:ch17-1775}下，存在序列 $\{\mu_j\}$（$\mu_j\to0$，$j\to\infty$），使得 $\{(u_{\mu_j},v_{\mu_j})\}$ 在上半平面上几乎处处收敛于\eqref{align:ch17-1771}的一个解 $(\bar u,\bar v)$，其初值为 $(\bar u(x,0),\bar v(x,0))=(u_0(x),v_0(x))$（$-\infty<x<\infty$）。
\end{theorem}

假设\eqref{align:ch17-1775}与特殊人工黏性\eqref{align:ch17-1776}的使用，在\cref{theorem:ch17-弹性存在性}的证明中是必不可少的，因为它们对于确立近似解的一致 $L^\infty$ 界似乎不可或缺。同时，有趣的是：人们是否可通过令黏弹性系统\eqref{align:ch08-弹性黏性}或至少模型系统
\begin{align}
\begin{cases}
\partial_t u-\partial_x v=0\\
\partial_t v-\partial_x\sigma(u)=\mu\,\partial_x^2v,
\end{cases}
\label{align:ch17-17722}
\end{align}
中 $\mu\to0$ 来构造\eqref{align:ch17-1771}的解，而该系统与上述模型系统相近。

尽管对\eqref{align:ch17-17722}的解我们没有一致 $L^\infty$ 估计（因为该系统相对于\eqref{align:ch17-1771}的所有凸熵并不耗散），但我们仍有若干 $L^p$ 型估计，其中最显著者是由物理熵--熵流对\eqref{align:ch07-p系统熵}诱导的“能量不等式”。
因此自然要问：补偿紧致性方法能否连同这类估计一起使用。当然，这将迫使我们放弃 $L^\infty$，并在 $L^p$ 框架下考虑 Young 测度——这一可能性在17.1节已提出。结果表明，该方法对眼前问题有效，尽管以精巧分析为代价，因此这里只记录结论。证明见\cref{section:ch17-notes}所引文献。

\begin{theorem}[等熵弹性的 $L^p$ 极限]\label{theorem:ch17-弹性Lp极限}
考虑系统\eqref{align:ch17-17722}，其中 (a) 对一切 $-\infty<u<\infty$，$\sigma'(u)\geqslant\sigma_0>0$；(b) $\sigma''$ 至多在 $(-\infty,\infty)$ 上一点处为零；(c) 当 $|u|\to\infty$ 时，$\sigma'(u)$ 增长如 $|u|^\alpha$（某 $\alpha\geqslant0$）；(d) 当 $|u|\to\infty$ 时，$\sigma''(u)$ 与 $\sigma'''(u)$ 增长不快于 $|u|^{\alpha-1}$。设 $(u_\mu,v_\mu)$ 是 Cauchy 问题\eqref{align:ch17-17722}、\eqref{align:ch17-1777}的解，其中 $\{(u_{0\mu},v_{0\mu})\}$ 是 $W^{1,2}(-\infty,\infty)$ 中具有一致有界总能量
\begin{align}
\int_{-\infty}^\infty\bigl[\tfrac{1}{2}v_{0\mu}^2(x)+e(u_{0\mu})\bigr]\,dx\leqslant a
\label{align:ch17-17723}
\end{align}
的函数，其振荡相对平缓，
\begin{align}
\mu\int_{-\infty}^\infty[\partial_x v_{0\mu}(x)]^2\,dx\longrightarrow0,\qquad\mu\to0,
\label{align:ch17-17724}
\end{align}
且当 $\mu\to0$ 时 $u_{0\mu}\to u_0$、$v_{0\mu}\to v_0$ 在分布意义下成立。则存在序列 $\{\mu_j\}$（$\mu_j\to0$，$j\to\infty$），使得 $\{(u_{\mu_j},v_{\mu_j})\}$ 对任意 $1<p<2$ 在 $L^p_{\mathrm{loc}}$ 中收敛于\eqref{align:ch17-1771}的一个解 $(u,v)$，其初值为 $(u_0,v_0)$（在 $(-\infty,\infty)$ 上）。
\end{theorem}

\end{document}
